\section*{Bloque C. Casos con implicancias regulatorias y sector saneamiento}

\noindent\textit{Preguntas 66 a 90. Evalúan la aplicación de la teoría de competencia imperfecta a subastas de Demsetz, licitaciones colusorias (bid rigging) en obras de EPS, facilidades esenciales, acceso de terceros, clemencia de Indecopi y regulación de servicios de saneamiento.}

\q{El mecanismo de 'Competencia por el Mercado' (Demsetz Auction) aplicado a la concesión de una Planta de Tratamiento de Agua Potable (PTAP) busca mitigar el poder de monopolio mediante:}
\begin{opts}
    \item subastar el contrato al postor que ofrezca la menor tarifa con calidad prefijada.
    \item la duplicación física de tuberías matrices para fomentar competencia en instalaciones.
    \item la eliminación de toda supervisión regulatoria posterior a la firma contractual.
    \item la fijación discrecional de precios máximos directamente por el ministerio sectorial.
\end{opts}

\q{En las licitaciones públicas convocadas por empresas prestadoras (EPS) para obras de ampliación de redes de alcantarillado, la práctica colusoria denominada 'posturas de resguardo' (cover bidding) consiste en:}
\begin{opts}
    \item solicitar la nulidad administrativa del proceso mediante apelaciones sistemáticas.
    \item presentar garantías financieras falsificadas ante el comité de adquisiciones.
    \item ofrecer descuentos por pronto pago fuera de los términos de las bases integradas.
    \item presentar ofertas altas o defectuosas para asegurar el triunfo del postor acordado.
\end{opts}

\q{Si los únicos tres fabricantes de insumos químicos para potabilización de agua (cloro gas y sulfato de aluminio) en el Perú acuerdan repartirse geográficamente a las EPS clientes (norte, centro y sur), esta conducta constituye:}
\begin{opts}
    \item una restricción vertical justificada en eficiencias logísticas de distribución.
    \item una discriminación de precios de segundo grado orientada al bienestar social.
    \item un monopolio bilateral eficiente acorde con la negociación directa de Coase.
    \item un cártel horizontal de reparto de mercado, infracción absoluta anticompetitiva.
\end{opts}

\q{Bajo la doctrina de las 'Facilidades Esenciales' (Essential Facilities Doctrine), una EPS integrada que opera una red matriz de distribución de agua comete abuso de posición de dominio si:}
\begin{opts}
    \item niega injustificadamente acceso a una red no duplicable pese a tener capacidad ociosa.
    \item aplica la estructura tarifaria regulada aprobada en el estudio tarifario de SUNASS.
    \item solicita autorización de vertimiento hídrico ante la Autoridad Nacional del Agua.
    \item exige cumplimiento de normas técnicas constructivas del Reglamento de Edificaciones.
\end{opts}

\q{Una EPS niega la factibilidad de servicios de agua potable a un nuevo desarrollo inmobiliario a pesar de tener capacidad ociosa en la tubería matriz vecina, exigiendo pagos extrarregulatorios. Esta conducta configura según Indecopi y SUNASS:}
\begin{opts}
    \item un ejercicio legítimo del derecho de propiedad exclusiva de la empresa municipal.
    \item una práctica lícita de discriminación de precios de primer grado al desarrollador.
    \item una externalidad positiva de red que beneficia a los usuarios residenciales.
    \item un abuso de posición de dominio por negativa injustificada de acceso a la red.
\end{opts}

\q{En las zonas periféricas no atendidas por la red pública de EPS, el mercado informal de camiones cisterna que distribuyen agua potable suele operar bajo una estructura de:}
\begin{opts}
    \item competencia perfecta con precio equivalente al costo marginal de extracción.
    \item monopsonio estatal regulado directamente por la autoridad sanitaria regional.
    \item oligopolio o competencia monopolística local con altos costos de transporte.
    \item monopolio natural puro con subaditividad global de costos medios decrecientes.
\end{opts}

\q{Bajo un esquema regulatorio de Precio Tope (Price Cap / RPI-X) aplicado a una EPS, si los contratistas de obras civiles de la empresa forman un cártel que infla los costos de inversión (CAPEX):}
\begin{opts}
    \item la empresa absorbe el sobrecosto de inversión sin impacto alguno en su rentabilidad.
    \item la EPS no traslada ese sobrecosto durante el quinquenio, reduciendo su rentabilidad.
    \item las tarifas residenciales se indexan mensualmente al costo facturado por la constructora.
    \item el factor de productividad X se reajusta inmediatamente a cero en el mes siguiente.
\end{opts}

\q{El Decreto Legislativo N° 1280 (Ley del Servicio Universal de Agua Potable y Saneamiento) promueve la integración de prestadores a nivel provincial buscando:}
\begin{opts}
    \item privatizar obligatoriamente la titularidad accionaria de todas las EPS municipales.
    \item fomentar la fragmentación distrital de operadores para competir en precios.
    \item aprovechar economías de escala y densidad frente a prestadores atomizados.
    \item sustituir la fiscalización tarifaria de SUNASS y las auditorías de gestión.
\end{opts}

\q{En el suministro de agua desalinizada para uso industrial o minero donde existen dos plantas independientes con límites estrictos de capacidad instalada, el modelo microeconómico más adecuado para describir la competencia es:}
\begin{opts}
    \item competencia perfecta de largo plazo con libre entrada instantánea sin costos hundidos.
    \item Bertrand-Edgeworth con restricciones de capacidad y precios sobre el costo marginal.
    \item Bertrand clásico con bienes homogéneos, costos simétricos y capacidad infinita.
    \item monopolio bilateral puro resuelto mediante negociación cooperativa simétrica.
\end{opts}

\q{Cuando una EPS adquiere micromedidores de agua mediante subasta pública, el riesgo de que los proveedores ofrezcan medidores con submedición intencional para ganar con el menor precio corresponde a un problema de:}
\begin{opts}
    \item selección adversa por asimetría de información sobre calidad técnica (lemons).
    \item riesgo moral ex-post posterior a la aprobación del pliego tarifario quinquenal.
    \item monopolio natural tecnológico en la fabricación global de caudalímetros.
    \item externalidad de red positiva generada hacia los grandes usuarios comerciales.
\end{opts}

\q{El 'Programa de Clemencia' (Leniency Program) administrado por Indecopi en la detección de cárteles de proveedores de obras de saneamiento consiste en:}
\begin{opts}
    \item otorgar subsidios estatales de rescate a constructoras con problemas de liquidez.
    \item autorizar acuerdos horizontales de precios si financian infraestructura rural.
    \item condonar las deudas tributarias acumuladas de las empresas prestadoras municipales.
    \item exonerar o reducir multas al primer delator que aporte pruebas clave del cártel.
\end{opts}

\q{La utilización de listas de precios de referencia públicamente conocidas o la coincidencia reiterada en cotizaciones con decimales idénticos entre postores en subastas de tuberías de saneamiento suele ser indicio de:}
\begin{opts}
    \item colusión tácita coordinada mediante puntos focales (Schelling) y señales indirectas.
    \item eficiencia productiva competitiva que maximiza el excedente neto del consumidor.
    \item rivalidad intensa en precios con información perfecta y transparente en subastas.
    \item comportamiento estrictamente no estratégico bajo costos marginales idénticos.
\end{opts}

\q{Si una empresa fabricante dominante de tuberías de PVC adquiere a la principal empresa contratista de obras de agua potable, las autoridades de competencia analizan el riesgo de:}
\begin{opts}
    \item cierre vertical de mercado (foreclosure) discriminando o excluyendo a contratistas rivales.
    \item colusión horizontal automática entre proveedores independientes de cloro líquido.
    \item aumento inmediato de la elasticidad precio en la demanda final de agua potable.
    \item surgimiento de un monopolio natural subaditivo en la producción básica de cemento.
\end{opts}

\q{Para determinar si las fuentes subterráneas privadas (pozos) pertenecen al mismo 'mercado relevante' que el agua potable por red de la EPS para grandes industrias, Indecopi aplica:}
\begin{opts}
    \item la regla de la ventaja comparativa ricardiana aplicada a recursos naturales.
    \item la prueba estadística de dominancia estocástica de primer orden de Bain.
    \item el Test del Monopolista Hipotético (SSNIP: aumento pequeño pero significativo).
    \item el modelo de equilibrio general competitivo intertemporal de Arrow-Debreu.
\end{opts}

\q{En el mercado de laboratorios de control de calidad de agua potable certificados por Inacal, operan dos laboratorios principales que compiten tanto en precios como en plazos de entrega. La interdependencia estratégica sugiere que:}
\begin{opts}
    \item los resultados fisicoquímicos perderán validez probatoria ante el regulador.
    \item fijarán tarifas idénticas al costo marginal aun con servicios diferenciados.
    \item sus precios actúan como complementos estratégicos con reacción positiva mutua.
    \item el laboratorio de menor escala cerrará automáticamente en el corto plazo.
\end{opts}

\q{La Regla de Paridad Eficiente de Componentes (ECPR / Baumol-Willig) para fijar el peaje de acceso que una EPS debe cobrar a un tercero por usar su red de transporte establece que el peaje debe ser igual a:}
\begin{opts}
    \item costo contable histórico promedio de amortización lineal de la tubería matriz.
    \item cero soles, para maximizar la competencia minorista sin considerar costos fijos.
    \item tarifa media final del servicio urbano dividida entre la densidad de usuarios.
    \item costo marginal directo del acceso más costo de oportunidad por beneficios cedidos.
\end{opts}

\q{Los subsidios cruzados en las tarifas de agua (donde usuarios comerciales e industriales pagan tarifas superiores al costo medio para subsidiar a usuarios residenciales vulnerables) pueden generar la siguiente distorsión estratégica:}
\begin{opts}
    \item anular por completo el consumo básico de los hogares en situación de pobreza.
    \item incentivar a grandes usuarios a desconectarse (bypass ineficiente) hacia pozos.
    \item eliminar la condición de monopolio natural en las redes de distribución local.
    \item generar un superávit fiscal automático transferible al Tesoro Público central.
\end{opts}

\q{Según el modelo de colusión con demanda fluctuante de Rotemberg y Saloner (1986), los acuerdos colusorios entre distribuidores de agua o insumos son más propensos a romperse (desatar guerras de precios) durante:}
\begin{opts}
    \item esquemas regulatorios estrictos de tasa de retorno sobre el capital invertido.
    \item periodos recesivos profundos con contracción prolongada de la demanda agregada.
    \item picos estacionales de demanda alta, cuando la ganancia de desviarse es máxima.
    \item contextos estacionarios donde el factor de descuento intertemporal iguala a 1.
\end{opts}

\q{Cuando una EPS necesita adquirir terrenos comunales para construir una Planta de Tratamiento de Aguas Residuales (PTAR) y no existen opciones técnicas alternativas de localización, la relación económica entre la EPS y la comunidad campesina se caracteriza como:}
\begin{opts}
    \item competencia perfecta bilateral con atomicidad de ofertantes y demandantes.
    \item mercado contestable de libre entrada con rendimientos crecientes a escala.
    \item monopolio bilateral donde el precio depende del poder relativo de negociación.
    \item duopolio de Cournot con costos de transporte lineales perfectamente simétricos.
\end{opts}

\q{En los contratos de Asociación Público-Privada (APP) cofinanciadas para obras de saneamiento en el Perú, el riesgo de que los postores incurran en la 'maldición del ganador' (Winner's Curse) implica que:}
\begin{opts}
    \item la tarifa regulada al usuario se duplica automáticamente tras suscribirse el contrato.
    \item el ganador subestimó costos u optimizó en exceso la demanda, arriesgando quiebra o adendas.
    \item la empresa adjudicataria resulta invariablemente la de mayor solvencia tecnológica.
    \item el postor derrotado en la subasta debe indemnizar monetariamente al Estado concedente.
\end{opts}

\q{Una empresa titular de una concesión de tratamiento de aguas residuales industriales instala deliberadamente capacidad de tratamiento ociosa y sobredimensionada. Estratégicamente, esta inversión en capital irreversible sirve para:}
\begin{opts}
    \item comprometer una respuesta agresiva que disuada creíblemente la entrada (Dixit).
    \item elevar el costo medio regulatorio para sustentar una reducción tarifaria futura.
    \item acreditar de manera automática las metas quinquenales de calidad ante SUNASS.
    \item convertir la función de costos totales en estrictamente lineal sin costos fijos.
\end{opts}

\q{Los altos costos de cambio (switching costs) que enfrenta un cliente no doméstico para desconectarse de la red de agua de la EPS e implementar su propio sistema de reúno o pozo:}
\begin{opts}
    \item garantizan que la elasticidad precio de la demanda industrial sea infinita.
    \item otorgan poder de mercado para cobrar sobre costos directos sin perder al cliente.
    \item eliminan toda barrera técnica para el ingreso de operadores descentralizados.
    \item intensifican la competencia directa en el segmento de clientes no domésticos.
\end{opts}

\q{La sobreexplotación desmedida de un acuífero subterráneo compartido por múltiples agroindustrias y la EPS local (ej. acuífero de Ica) representa un resultado del tipo:}
\begin{opts}
    \item monopolio natural sostenible sujeto a precios óptimos de Ramsey-Boiteux.
    \item juego cooperativo estable con transferencias laterales plenamente ejecutables.
    \item equilibrio de Nash eficiente en el sentido de Pareto con conservación hídrica.
    \item Dilema del Prisionero n-personal (Tragedia de los Comunes) que agota el recurso.
\end{opts}

\q{La frontera de competencias entre la SUNASS y el Indecopi en el sector saneamiento establece que:}
\begin{opts}
    \item SUNASS fiscaliza acuerdos horizontales en insumos químicos e Indecopi aprueba los pliegos tarifarios.
    \item ambas entidades ejercen potestad tarifaria simultánea e imponen sanciones duplicadas por igual hecho.
    \item SUNASS regula tarifas, calidad y acceso de EPS; Indecopi sanciona cárteles y abusos de dominio.
    \item ninguna autoridad sectorial ejerce supervisión económica sobre empresas prestadoras de capital estatal.
\end{opts}

\q{Una EPS condiciona la aprobación de la conexión de agua potable para una nueva urbanización a que los desarrolladores compren obligatoriamente los medidores y tuberías a una filial comercial de la EPS. Esta práctica se clasifica como:}
\begin{opts}
    \item un estándar de calidad técnica equivalente a la homologación metrológica oficial.
    \item un subsidio cruzado normado explícitamente bajo el Decreto Legislativo N° 1280.
    \item una discriminación de precios de tercer grado segmentada por uso constructivo.
    \item una venta atada (tying) que traslada poder monopólico al mercado de materiales.
\end{opts}

